\section{Otomata dan Daya Ekspresif}
\label{sec:ch04-otomata}

Tabel \tabref{tab:ch04-ekivalensi} merangkum nama mesin pengenal untuk
setiap kelas pada hirarki Chomsky. Nama-nama ini dipakai pada tabel
pemetaan Sub-CPMK 1.4; rincian konstruksi dan simulasinya dibahas pada
Bab VI.

\begin{table}[H]
\centering
\caption{Hubungan tata bahasa, bahasa, dan mesin}
\label{tab:ch04-ekivalensi}
\begin{tabularx}{\textwidth}{@{}l X l@{}}
\toprule
\textbf{Tata bahasa} & \textbf{Bahasa} & \textbf{Automata/Mesin} \\
\midrule
Regular & Regular & DFA/NFA \\
Bebas konteks & Bebas konteks & PDA \\
Sensitif konteks & Sensitif konteks & Otomata terbatas linear \\
Tanpa batas & Rekursif terbilang & Mesin Turing \\
\bottomrule
\end{tabularx}
\end{table}

\subsection{Nama yang Perlu Dibedakan}

\begin{description}[leftmargin=*]
  \item[DFA/NFA] Otomata hingga deterministik dan nondeterministik.
        Keduanya mengenali bahasa regular dan menjadi dasar scanner.
  \item[PDA] Otomata pushdown, yaitu otomata hingga yang dilengkapi
        stack. PDA mengenali bahasa bebas konteks dan menjadi model
        parser.
  \item[Otomata terbatas linear] Mesin Turing yang pitanya dibatasi
        sepanjang string masukan. Mesin ini mengenali bahasa sensitif
        konteks.
  \item[Mesin Turing] Model komputasi dengan pita tak terbatas. Mesin
        ini mengenali bahasa rekursif terbilang dan menjadi batas atas
        daya komputasi.
\end{description}

Ekivalensi ini bukan sekadar hasil teoretis. Ia menjadi dasar
pemilihan alat: ekspresi reguler dan \eng{lexer generator} untuk
token, CFG dan \eng{parser generator} untuk struktur, serta analisis
atribut untuk semantik. Memaksakan alat yang lebih kuat dari yang
diperlukan biasanya menghasilkan solusi yang lebih lambat dan lebih
sulit dipelihara.
